\documentclass[preview]{standalone}

\usepackage{amsmath}
\usepackage{amssymb}
\usepackage{stellar}
\usepackage{definitions}

\begin{document}

\id{psicologia-personalita-psicodinamica}
\genpage

\section{Teorie psicodinamiche}

\begin{snippetdefinition}{teorie-psicodinamiche-definition}{Teorie psicodinamiche}
    Le \emph{teorie psicodinamiche} della personalità condividono l'assunto
    secondo cui potenti forze interiori plasmano la personalità e motivano il
    comportamento.
\end{snippetdefinition}

\begin{snippet}{freud-personalita}
    Le teorie psicodinamiche prendono avvio da Sigmund Freud. La teoria
    freudiana cerca di spiegare le origini e lo sviluppo della personalità, la
    natura della mente, gli aspetti atipici della personalità e il modo in cui
    essa può essere modificata attraverso la terapia.
\end{snippet}

\begin{snippetdefinition}{libido-definition}{Libido}
    La \emph{libido} è, secondo Freud, la fonte di energia degli impulsi
    sessuali.
\end{snippetdefinition}

\begin{snippetdefinition}{destrudo-definition}{Destrudo}
    La \emph{destrudo} è la componente energetica associata alla pulsione
    aggressiva.
\end{snippetdefinition}

\begin{snippet}{sviluppo-psicosessuale}
    Freud riteneva che la pulsione sessuale agisse fin dalla nascita e che si
    organizzasse in fasi dello sviluppo psicosessuale: orale, anale, fallica, di
    latenza e genitale. I conflitti delle prime fasi di vita possono influenzare
    i comportamenti successivi.
\end{snippet}

\begin{snippetdefinition}{fissazione-psicodinamica-definition}{Fissazione}
    La \emph{fissazione} è l'incapacità di accedere alla fase successiva dello
    sviluppo psicosessuale, causata da eccessiva gratificazione o eccessiva
    frustrazione in una delle prime fasi.
\end{snippetdefinition}

\begin{snippetdefinition}{determinismo-psichico-definition}{Determinismo psichico}
    Il \emph{determinismo psichico} è l'assunto secondo cui tutti gli eventi
    della vita psichica sono determinati da esperienze precedenti e da episodi
    significativi, consci o inconsci.
\end{snippetdefinition}

\begin{snippet}{manifesto-latente}
    Per Freud, il comportamento può avere un \emph{contenuto manifesto}, cioè
    ciò che si dice, si fa e si percepisce consapevolmente, e un
    \emph{contenuto latente}, cioè il significato nascosto e non immediatamente
    consapevole.
\end{snippet}

\begin{snippetdefinition}{es-definition}{Es}
    L'\emph{Es} è la struttura primitiva e inconscia della personalità. Opera in
    modo irrazionale, agisce d'impulso ed esige gratificazione immediata secondo
    il principio del piacere, senza considerare realtà, norme sociali o moralità.
\end{snippetdefinition}

\begin{snippetdefinition}{io-definition}{Io}
    L'\emph{Io} è l'aspetto della personalità coinvolto nelle attività di
    autoconservazione e nel dirigere le pulsioni verso sbocchi appropriati.
    Regola il rapporto tra individuo e ambiente e opera secondo il principio di
    realtà.
\end{snippetdefinition}

\begin{snippetdefinition}{super-io-definition}{Super-Io}
    Il \emph{Super-Io} è l'istanza della psiche che comprende norme sociali e
    morali e aspirazioni ideali. Corrisponde in parte alla coscienza morale ed è
    spesso in tensione con l'Es.
\end{snippetdefinition}

\begin{snippet}{mediazione-io}
    L'Io media tra Es e Super-Io, sostituendo alla ricerca immediata del piacere
    scelte basate sulla ragione e sulla realtà. Quando si trova davanti a stati
    di angoscia, mette in atto meccanismi di difesa.
\end{snippet}

\begin{snippetdefinition}{rimozione-definition}{Rimozione}
    La \emph{rimozione} è il meccanismo di difesa attraverso cui pensieri,
    sentimenti o ricordi percepiti come minacciosi vengono esclusi dalla
    coscienza.
\end{snippetdefinition}

\begin{snippet}{valutazione-freud}
    Le idee di Freud hanno avuto grande impatto sull'interpretazione della
    personalità ordinaria e patologica, ma sono state criticate perché molti
    concetti risultano vaghi e difficili da definire o verificare empiricamente.
\end{snippet}

\begin{snippet}{post-psicoanalisi}
    Gli studiosi successivi a Freud attribuirono maggiore importanza alle
    funzioni dell'Io, alle variabili sociali, alla cultura, alla famiglia, ai
    pari e allo sviluppo della personalità oltre l'infanzia, riducendo
    l'enfasi sugli impulsi sessuali generali.
\end{snippet}

\begin{snippetdefinition}{inconscio-personale-definition}{Inconscio personale}
    L'\emph{inconscio personale}, secondo Jung, è formato da contenuti
    dimenticati o rimossi dall'individuo.
\end{snippetdefinition}

\begin{snippetdefinition}{inconscio-collettivo-definition}{Inconscio collettivo}
    L'\emph{inconscio collettivo}, secondo Jung, è una dimensione inconscia
    ereditata e condivisa dall'intera specie umana.
\end{snippetdefinition}

\begin{snippetdefinition}{archetipo-definition}{Archetipo}
    Un \emph{archetipo} è un'immagine universale e innata, collocata
    nell'inconscio collettivo, che orienta il comportamento umano.
\end{snippetdefinition}

\begin{snippetdefinition}{psicologia-analitica-definition}{Psicologia analitica}
    La \emph{psicologia analitica} è la prospettiva di Jung secondo cui la
    personalità è una costellazione di forze interiori che si compensano in
    equilibrio dinamico, con enfasi sul bisogno di realizzazione del sé.
\end{snippetdefinition}

\end{document}
